\documentclass{beamer}

\usepackage{beamerthemeSingapore}
\usepackage[brazil]{babel}
\usepackage[T1]{fontenc}
\usepackage[utf8]{inputenc}
\usepackage{graphicx}
\usepackage{url}
\usepackage{xcolor}
\usepackage{booktabs}
\usepackage{ragged2e}
\usepackage{etoolbox}
\usepackage{minted}
\usepackage{tikz}
\usetikzlibrary{positioning, arrows.meta, shapes.geometric}

\newcommand{\tcb}[1]{\textcolor{blue}{#1}}
\newcommand{\tcr}[1]{\textcolor{red}{#1}}
\newcommand{\tcg}[1]{\textcolor{green!50!black}{#1}}
\newcommand{\ddash}{-\mbox{}-}

\setminted{
  fontsize=\footnotesize,
  breaklines=true,
  breakanywhere=true,
  frame=lines,
  framesep=2mm,
  baselinestretch=1.05,
  tabsize=2
}

\apptocmd{\frame}{}{\justifying}{}

\setbeamertemplate{navigation symbols}{}

\title{05 - CSS3: layout e responsividade}
\author{Prof. Alex Kutzke}
\date{Setembro de 2026}

\begin{document}

% ==========================================
\begin{frame}
	\begin{center}
		\large{Setor de Educação Profissional e Tecnológica}\\
		Tecnologia em Análise e Desenvolvimento de Sistemas\\
		\vspace{1cm}
		\small{DS122 - Desenvolvimento de Aplicações Web 1}
	\end{center}
	\titlepage
\end{frame}

% ==========================================
\begin{frame}{Objetivos da aula}
	Ao final desta aula você deve ser capaz de:
	\begin{enumerate}
		\item Dizer o que muda nos filhos quando o pai recebe \texttt{display: flex}
		      ou \texttt{display: grid};
		\item Montar uma barra de navegação horizontal com Flexbox;
		\item Montar uma grade de cartões com Grid e calcular quantas colunas ela
		      terá em uma largura dada;
		\item Escolher entre Flexbox e Grid justificando pela dimensão do problema;
		\item Escrever \emph{media queries} em \emph{mobile-first};
		\item Verificar a página a 320px e corrigir a rolagem horizontal;
		\item Ler e alterar HTML que usa a grade e as utilitárias do Bootstrap 5;
		\item Usar o inspetor de Flexbox e de Grid do navegador.
	\end{enumerate}
\end{frame}

% ==========================================
\begin{frame}{Roteiro}
	\tableofcontents
\end{frame}

% ==========================================
\section{Fluxo}
% ==========================================

\begin{frame}{Onde a página parou}
	O \texttt{catalogo.html} da tarefa passada tem borda, \texttt{padding} e cor
	de fundo em cada cartão.

	\vspace{0.3cm}
	E os cartões continuam empilhados, um debaixo do outro, ocupando a largura
	inteira.

	\vspace{0.5cm}
	Nada de errado aconteceu: isso é o \textbf{fluxo normal}.
\end{frame}

\begin{frame}[fragile]{Três declarações no inspetor}
	Com \texttt{F12} aberto, na regra do elemento que contém os cartões:

	\begin{minted}{css}
display: flex;                        /* mesma linha */
gap: 20px;                            /* afastados */

display: grid;                        /* três colunas */
grid-template-columns: repeat(3, 1fr);
  \end{minted}

	\vspace{0.2cm}
	Depois, \texttt{Ctrl+Shift+M} liga o modo responsivo. Arraste até uma largura
	de celular e observe o texto transbordar.
\end{frame}

\begin{frame}{O que apareceu na tela}
	\begin{itemize}
		\item o elemento onde você escreveu o \texttt{display} passou a mandar na
		      posição dos filhos;
		\item ao lado do seletor surgiu um distintivo \texttt{flex} ou
		      \texttt{grid}, que liga um desenho por cima da página;
		\item a largura da janela virou uma variável do problema.
	\end{itemize}

	\vspace{0.5cm}
	A aula nomeia essas três coisas, e volta ao inspetor no fim.
\end{frame}

\begin{frame}{Por que os cartões empilham}
	Fluxo normal: caixas \texttt{block} uma abaixo da outra, na ordem do
	documento, ocupando toda a largura; caixas \texttt{inline} lado a lado,
	quebrando quando falta espaço.

	\vspace{0.3cm}
	\texttt{<article>}, \texttt{<section>} e \texttt{<div>} são \texttt{block}.

	\vspace{0.5cm}
	Até meados dos anos 2010, pôr dois blocos lado a lado exigia \texttt{float},
	criada para texto contornar imagem, ou tabelas de diagramação. As duas
	soluções eram frágeis e exigiam remendos.
\end{frame}

\begin{frame}[fragile]{Contêiner e item}
	\begin{minted}{html}
<div class="grade">          <!-- contêiner -->
  <article class="produto">  <!-- item -->
    <h3>Café em grão</h3>    <!-- não é item -->
    <p>Torra média.</p>
  </article>
  <article class="produto">  <!-- item -->
  </article>
</div>
  \end{minted}

	\vspace{0.2cm}
	O \textbf{contêiner} recebe \texttt{display: flex} ou \texttt{display: grid}.
	Os \textbf{itens} são os \textbf{filhos diretos} dele. Netos voltam ao fluxo
	normal dentro do próprio pai.
\end{frame}

\begin{frame}{Duas famílias de propriedades}
	Propriedades que se escrevem \textbf{no contêiner} e propriedades que se
	escrevem \textbf{no item}.

	\vspace{0.4cm}
	\begin{block}{O erro mais comum}
		\texttt{justify-content} dentro de \texttt{.produto} não faz nada, e o
		navegador não avisa. CSS inválido ou inútil é descartado em silêncio.
	\end{block}

	\vspace{0.4cm}
	Quando uma propriedade de layout parecer ignorada, a primeira pergunta é se
	ela está no elemento certo.
\end{frame}

% ==========================================
\section{Flexbox}
% ==========================================

\begin{frame}{Uma direção, e a sobra repartida}
	O Flexbox distribui itens ao longo de \textbf{uma} direção, e reparte entre
	eles o espaço que sobra ou que falta.

	\vspace{0.5cm}
	\begin{center}
		\begin{tabular}{ll}
			\toprule
			\textbf{Propriedade}        & \textbf{Atua no eixo} \\
			\midrule
			\texttt{justify-content}    & principal             \\
			\texttt{align-items}        & transversal           \\
			\bottomrule
		\end{tabular}
	\end{center}

	\vspace{0.4cm}
	O eixo principal é definido por \texttt{flex-direction} (padrão
	\texttt{row}). Com \texttt{column}, as duas propriedades trocam de sentido.
	Guarde a relação pelo eixo, e não por horizontal e vertical.
\end{frame}

\begin{frame}[fragile]{Propriedades do contêiner}
	\begin{minted}{css}
.barra {
  display: flex;
  flex-direction: row;            /* row | column | ... */
  justify-content: space-between; /* eixo principal */
  align-items: center;            /* eixo transversal */
  flex-wrap: wrap;                /* permite quebrar linha */
  gap: 1rem;                      /* espaço entre itens */
}
  \end{minted}
\end{frame}

\begin{frame}{Os valores que aparecem na prática}
	\begin{center}
		\begin{tabular}{ll}
			\toprule
			\multicolumn{2}{c}{\textbf{justify-content}}                            \\
			\midrule
			\texttt{flex-start}    & itens no começo do eixo (padrão)               \\
			\texttt{center}        & itens no meio                                  \\
			\texttt{space-between} & pontas nas bordas, sobra dividida entre eles   \\
			\texttt{space-evenly}  & espaços iguais, inclusive nas pontas           \\
			\midrule
			\multicolumn{2}{c}{\textbf{align-items}}                                \\
			\midrule
			\texttt{stretch}       & esticam até a altura do mais alto (padrão)     \\
			\texttt{center}        & centralizados na altura                        \\
			\texttt{flex-start}    & alinhados pelo topo                            \\
			\texttt{baseline}      & alinhados pela linha de base do texto          \\
			\bottomrule
		\end{tabular}
	\end{center}
\end{frame}

\begin{frame}{Duas armadilhas}
	\begin{block}{\texttt{flex-wrap: nowrap} é o padrão}
		Sem \texttt{wrap}, os itens nunca quebram para a linha de baixo: eles
		encolhem até caber, e continuam encolhendo depois que o conteúdo já não
		cabe mais. Em cartões de produto, quase sempre você quer \texttt{wrap}.
	\end{block}

	\vspace{0.4cm}
	\begin{block}{\texttt{gap} não é \texttt{margin}}
		\texttt{gap} vale só \emph{entre} os itens, e não sobra espaço nas pontas,
		que é o que a margem faria. Funciona em Flexbox e em Grid.
	\end{block}
\end{frame}

\begin{frame}[fragile]{Propriedades do item}
	\begin{minted}{css}
.item {
  flex-grow: 1;       /* fatia da sobra (padrão 0) */
  flex-shrink: 1;     /* aceita encolher (padrão 1) */
  flex-basis: auto;   /* tamanho de partida */
  align-self: center; /* alinhamento só deste item */
}
  \end{minted}

	\vspace{0.2cm}
	A forma abreviada resolve quase tudo:

	\begin{minted}{css}
.busca { flex: 1; }  /* = 1 1 0%: ocupe toda a sobra */
.logo  { flex: 0; }  /* tamanho do conteúdo */
  \end{minted}
\end{frame}

\begin{frame}[fragile]{Uma barra com campo elástico}
	\begin{minted}{css}
.barra { display: flex; align-items: center; gap: 1rem; }

.barra .logo   { flex: 0 0 auto; } /* não encolhe */
.barra .busca  { flex: 1 1 auto; } /* absorve a sobra */
.barra .entrar { flex: 0 0 auto; }
  \end{minted}

	\vspace{0.3cm}
	Logotipo à esquerda, busca no meio ocupando o que sobrar, botão à direita.
\end{frame}

\begin{frame}[fragile]{O padrão da barra de navegação}
	\begin{minted}{css}
header {
  display: flex;
  justify-content: space-between;
  align-items: center;
  padding: 1rem;
}

.menu ul {
  display: flex;
  gap: 1.5rem;
  list-style: none;  /* tira o marcador */
  margin: 0;
  padding: 0;        /* tira o recuo padrão da lista */
}
  \end{minted}

	\vspace{0.2cm}
	Só as duas primeiras declarações do \texttt{ul} mudam o layout. As outras três
	a lista precisaria de qualquer forma.
\end{frame}

\begin{frame}[fragile]{\texttt{order} e a ordem de leitura}
	\begin{minted}{css}
.destaque { order: -1; }  /* aparece antes dos demais */
  \end{minted}

	\vspace{0.3cm}
	A ordem \textbf{visual} muda. A ordem do \textbf{documento} não.

	\vspace{0.3cm}
	Quem navega com \texttt{Tab} e quem usa leitor de tela continua percorrendo os
	elementos na ordem do HTML. Se as duas ordens discordarem, a pessoa vê o foco
	pular de um canto a outro da tela.

	\vspace{0.3cm}
	WCAG 1.3.2 (sequência com significado) e 2.4.3 (ordem de foco). Quando a ordem
	certa for outra, mude o HTML.
\end{frame}

% ==========================================
\section{Grid}
% ==========================================

\begin{frame}[fragile]{Linhas e colunas ao mesmo tempo}
	\begin{minted}{css}
.grade {
  display: grid;
  grid-template-columns: 1fr 1fr 1fr; /* três colunas iguais */
  gap: 20px;
}
  \end{minted}

	\vspace{0.3cm}
	\texttt{fr} é uma fração do espaço livre do contêiner, descontados o
	\texttt{gap} e as colunas de tamanho fixo. Só existe dentro do Grid.

	\begin{minted}{css}
/* mesma coisa que 1fr 1fr 1fr */
grid-template-columns: repeat(3, 1fr);

/* lateral fixa de 250px, e o resto para o conteúdo */
grid-template-columns: 250px 1fr;
  \end{minted}
\end{frame}

\begin{frame}[fragile]{Um item que ocupa mais de uma célula}
	\begin{minted}{css}
.produto.destaque {
  grid-column: span 2;  /* ocupa duas colunas */
  grid-row: span 2;     /* e duas linhas */
}
  \end{minted}

	\vspace{0.4cm}
	É o recurso que o Flexbox não tem, e o motivo de o Grid existir.

	\vspace{0.4cm}
	\texttt{grid-template-rows} existe, mas raramente é necessária: as linhas que
	faltarem são criadas automaticamente, com a altura do conteúdo.
\end{frame}

\begin{frame}[fragile]{Áreas nomeadas: o esqueleto da página}
	\begin{minted}{css}
.pagina {
  display: grid;
  grid-template-columns: 1fr 250px;
  grid-template-areas:
    "cabecalho cabecalho"
    "conteudo  lateral"
    "rodape    rodape";
  gap: 1rem;
}

header { grid-area: cabecalho; }
main   { grid-area: conteudo; }
aside  { grid-area: lateral; }
footer { grid-area: rodape; }
  \end{minted}
\end{frame}

\begin{frame}[fragile]{O mesmo layout na tela estreita}
	Cada linha entre aspas é uma linha da grade. Nome repetido em células vizinhas
	estende a área. Um ponto marca célula vazia.

	\begin{minted}{css}
grid-template-columns: 1fr;
grid-template-areas:
  "cabecalho"
  "conteudo"
  "lateral"
  "rodape";
  \end{minted}

	\vspace{0.2cm}
	Reescrever o layout é redesenhar o quadro.
\end{frame}

\begin{frame}[fragile]{A grade que se adapta sem media query}
	\begin{minted}{css}
.grade {
  display: grid;
  grid-template-columns: repeat(auto-fit, minmax(240px, 1fr));
  gap: 20px;
}
  \end{minted}

	\vspace{0.3cm}
	Em português: crie quantas colunas couberem, com no mínimo 240px cada uma, e
	reparta a sobra igualmente entre elas.

	\vspace{0.3cm}
	Esta declaração resolve a grade de cartões do trabalho prático inteira.
\end{frame}

\begin{frame}{A conta, e a medição}
	Com \texttt{gap} de 20px, cabem \texttt{n} colunas enquanto

	\begin{center}
		\large $n \times 240 + (n-1) \times 20 \leq$ largura do contêiner
	\end{center}

	\vspace{0.2cm}
	Em 1000px: 4 colunas dariam $960 + 60 = 1020$, não cabe. 3 colunas dão
	$720 + 40 = 760$, cabe. Cada uma recebe $(1000 - 40)/3 = 320$px.

	\vspace{0.3cm}
	\begin{center}
		\begin{tabular}{lll}
			\toprule
			\textbf{Contêiner} & \textbf{Colunas} & \textbf{Cada uma} \\
			\midrule
			1000px             & 3                & 320px             \\
			800px              & 3                & 253.33px          \\
			600px              & 2                & 290px             \\
			500px              & 2                & 240px             \\
			400px              & 1                & 400px             \\
			\bottomrule
		\end{tabular}
	\end{center}
\end{frame}

\begin{frame}{\texttt{auto-fit} e \texttt{auto-fill}}
	\texttt{auto-fill} cria as colunas vazias em vez de descartá-las.

	\vspace{0.4cm}
	Com poucos itens e muito espaço:

	\begin{itemize}
		\item \texttt{auto-fill} deixa os cartões no tamanho mínimo e sobra vazio à
		      direita;
		\item \texttt{auto-fit} estica os cartões que existem.
	\end{itemize}

	\vspace{0.4cm}
	Para catálogo, \texttt{auto-fit} costuma ser o que se quer.
\end{frame}

\begin{frame}{Flexbox ou Grid}
	\begin{center}
		\begin{tabular}{p{4.2cm}p{4.8cm}}
			\toprule
			\textbf{Flexbox quando}                      & \textbf{Grid quando}                                  \\
			\midrule
			os itens ficam em uma linha ou em uma coluna & há linhas e colunas ao mesmo tempo                    \\
			\addlinespace
			o tamanho de cada item vem do conteúdo dele  & o desenho da grade vem primeiro                       \\
			\addlinespace
			distribuir a sobra entre poucos itens        & alinhar itens de linhas diferentes                    \\
			\bottomrule
		\end{tabular}
	\end{center}

	\vspace{0.4cm}
	Na página de catálogo: Grid no esqueleto e na grade de produtos; Flexbox no
	cabeçalho, no menu e dentro do cartão. Um contêiner Flex pode conter um Grid,
	e o contrário também.

	\vspace{0.2cm}
	Nos dois modelos, as margens verticais dos itens não colapsam.
\end{frame}

% ==========================================
\section{Responsividade}
% ==========================================

\begin{frame}[fragile]{A meta viewport}
	Sem esta linha no \texttt{<head>}, nada do que vem depois funciona no celular:

	\begin{minted}[fontsize=\scriptsize]{html}
<meta name="viewport" content="width=device-width, initial-scale=1">
  \end{minted}

	\vspace{0.3cm}
	Sem ela, o navegador do celular finge uma janela de cerca de \textbf{980px},
	desenha a página nela e reduz tudo por escala. Resultado: a página em
	miniatura, texto ilegível, e nenhuma \emph{media query} de largura pequena
	entra em ação.

	\vspace{0.3cm}
	\tcr{Não} acrescente \texttt{user-scalable=no} nem \texttt{maximum-scale=1}:
	elas impedem o zoom com os dedos, e quem tem baixa visão depende disso.
\end{frame}

\begin{frame}[fragile]{Media queries}
	\begin{minted}{css}
/* estilo base: telas pequenas */
.grade { grid-template-columns: 1fr; }

@media (min-width: 768px) {
  .grade { grid-template-columns: repeat(2, 1fr); }
}

@media (min-width: 1024px) {
  .grade { grid-template-columns: repeat(3, 1fr); }
}
  \end{minted}

	\vspace{0.2cm}
	Dentro do bloco vão regras CSS comuns, com os mesmos seletores.
\end{frame}

\begin{frame}{Media query não muda especificidade}
	As regras de dentro do bloco vencem as de fora por virem \textbf{depois} no
	arquivo, e não por serem mais específicas.

	\vspace{0.4cm}
	\begin{block}{Consequência prática}
		As \emph{media queries} vão ao \textbf{final} da folha. Escritas antes das
		regras que pretendem sobrescrever, elas perdem em silêncio.
	\end{block}
\end{frame}

\begin{frame}{Mobile-first}
	CSS base para a tela pequena, ampliado com \texttt{min-width}.

	\vspace{0.4cm}
	O caminho inverso, começar pelo desktop e cortar com \texttt{max-width},
	também funciona, e produz folhas maiores: o layout simples costuma ser o da
	tela estreita, e ele fica sendo o caso que não precisa de regra nenhuma.

	\vspace{0.4cm}
	Há também a razão de origem: a maior parte do acesso à web vem de telefone, e
	frequentemente de aparelho modesto com conexão limitada.
\end{frame}

\begin{frame}{O ponto de quebra vem do conteúdo}
	A tentação é copiar uma lista de larguras de aparelhos. Essa lista muda todo
	ano e nunca esteve completa.

	\vspace{0.4cm}
	O critério prático: arraste a janela do estreito para o largo e observe onde o
	layout fica ruim. Ali é o ponto de quebra.

	\vspace{0.4cm}
	Duas ou três quebras resolvem a maior parte das páginas desta disciplina.
\end{frame}

\begin{frame}[fragile]{Imagens e mídia fluidas}
	Uma imagem de 1200px de largura natural estoura qualquer coluna estreita e
	cria rolagem horizontal na página inteira.

	\begin{minted}{css}
img, video {
  max-width: 100%;
  height: auto;
}
  \end{minted}

	\vspace{0.2cm}
	\texttt{max-width: 100\%} impede a imagem de passar da largura do contêiner.
	\texttt{height: auto} mantém a proporção enquanto ela encolhe.
\end{frame}

\begin{frame}{O teste dos 320px}
	\begin{block}{WCAG 2.2, critério 1.4.10 (\emph{Reflow}, nível AA)}
		A página deve continuar utilizável, sem rolagem horizontal, em uma largura
		de \textbf{320 pixels CSS}.
	\end{block}

	\vspace{0.3cm}
	Equivale a um celular pequeno, e ao que sobra quando alguém amplia a página em
	400\% num monitor comum.

	\vspace{0.3cm}
	Modo responsivo, largura 320, rolar a página de cima a baixo. Se aparecer
	barra horizontal, os suspeitos de sempre são imagem sem \texttt{max-width},
	largura fixa em \texttt{px}, \texttt{white-space: nowrap} e tabela larga.
\end{frame}

\begin{frame}{Layout, acessibilidade e os ODS}
	\begin{itemize}
		\item grade que não se adapta obriga quem acessa pelo telefone a ampliar e
		      arrastar a tela a cada linha de texto;
		\item ordem visual em desacordo com a do documento desorienta quem navega
		      por teclado;
		\item página que só funciona em conexão rápida exclui quem não tem uma.
	\end{itemize}

	\vspace{0.4cm}
	\textbf{ODS 4} e \textbf{ODS 10}: o acesso à informação depende de a página
	funcionar no aparelho que a pessoa tem, e não no de quem a desenvolveu.
	\textbf{ODS 8}: interface que funciona no dispositivo de quem opera o sistema
	reduz retrabalho de quem trabalha com ela todo dia.
\end{frame}

% ==========================================
\section{Bootstrap}
% ==========================================

\begin{frame}[fragile]{Um framework CSS}
	Folha de estilo pronta, escrita por terceiros, usada pelas classes que ela
	define. O Bootstrap é o mais difundido, e é o que a ementa pede que você
	conheça.

	\begin{minted}[fontsize=\tiny]{html}
<link href="https://cdn.jsdelivr.net/npm/bootstrap@5.3.8/dist/css/bootstrap.min.css"
      rel="stylesheet"
      integrity="sha384-sRIl4kxILFvY47J16cr9ZwB07vP4J8+LH7qKQnuqkuIAvNWLzeN8tE5YBujZqJLB"
      crossorigin="anonymous">
  \end{minted}

	\vspace{0.2cm}
	\texttt{integrity} traz o resumo criptográfico do arquivo esperado. Se o que
	chegar do CDN não corresponder, o navegador recusa o arquivo.
\end{frame}

\begin{frame}{A grade de 12 colunas}
	\begin{center}
		\includegraphics[width=0.95\textwidth]{../../src/images/bootstrap/colunas.png}
	\end{center}

	\vspace{0.3cm}
	A largura é dividida em 12 colunas. Você declara quantas cada bloco ocupa, e a
	soma de uma linha deve dar 12.
\end{frame}

\begin{frame}[fragile]{Grade e pontos de quebra}
	\begin{minted}{html}
<div class="container">
  <div class="row">
    <div class="col-md-8">Conteúdo principal</div>
    <div class="col-md-4">Barra lateral</div>
  </div>
</div>
  \end{minted}

	\begin{center}
		\begin{tabular}{ll}
			\toprule
			\textbf{Sufixo} & \textbf{A partir de} \\
			\midrule
			(nenhum)        & qualquer largura     \\
			\texttt{sm}     & 576px                \\
			\texttt{md}     & 768px                \\
			\texttt{lg}     & 992px                \\
			\texttt{xl}     & 1200px               \\
			\texttt{xxl}    & 1400px               \\
			\bottomrule
		\end{tabular}
	\end{center}
\end{frame}

\begin{frame}[fragile]{Classes utilitárias}
	\begin{minted}[fontsize=\scriptsize]{html}
<div class="d-flex justify-content-between
            align-items-center p-3 mb-4 bg-light">
  \end{minted}

	\vspace{0.3cm}
	\texttt{d-flex} é \texttt{display: flex}; \texttt{p-3} é um \texttt{padding}
	da escala de espaçamento; \texttt{mb-4} é \texttt{margin-bottom};
	\texttt{bg-light} é uma cor de fundo.

	\vspace{0.3cm}
	Por baixo, \texttt{.row} é um contêiner Flexbox e \texttt{.col-md-8} define
	largura em porcentagem dentro de uma \emph{media query}. É o CSS desta aula,
	empacotado.
\end{frame}

\begin{frame}{O que isso custa, e a posição da disciplina}
	O \texttt{bootstrap.min.css} 5.3.8 tem \textbf{232.111 bytes}, cerca de
	227\,KB, que chegam como cerca de 30\,KB comprimidos. A página carrega a folha
	inteira, mesmo usando cinco classes dela. E o HTML acumula classes de
	apresentação, que é o que a separação entre estrutura e estilo pretendia
	evitar.

	\vspace{0.4cm}
	\begin{block}{Nesta disciplina}
		O Bootstrap não é ferramenta obrigatória, e o trabalho prático é feito com
		CSS escrito por você. Quem aprende \texttt{col-md-4} sem ter escrito uma
		\emph{media query} fica sem saber o que fazer quando a grade do framework
		não resolve o caso.

		Você vai encontrá-lo em código de terceiros e em estágio: o que se espera
		daqui é saber \textbf{ler}.
	\end{block}
\end{frame}

% ==========================================
\section{Prática}
% ==========================================

\begin{frame}{Roteiro da prática}
	Trabalhando sobre as páginas da tarefa de CSS:

	\begin{enumerate}
		\item meta viewport nas três páginas, e o menu na horizontal;
		\item a grade de produtos com \texttt{auto-fit} e \texttt{minmax};
		\item o cartão por dentro, com o preço alinhado na base;
		\item \emph{media queries} e o teste dos 320px.
	\end{enumerate}

	\vspace{0.4cm}
	O texto completo do roteiro está no material da aula, na parte prática.
\end{frame}

\begin{frame}[fragile]{O menu na horizontal}
	\begin{minted}{css}
header {
  display: flex;
  justify-content: space-between;
  align-items: center;
  gap: 1rem;
  padding: 1rem;
}

.menu ul {
  display: flex; gap: 1.5rem;
  list-style: none; margin: 0; padding: 0;
}
  \end{minted}

	\vspace{0.2cm}
	Dois experimentos, um de cada vez: trocar \texttt{space-between} por
	\texttt{center}; acrescentar \texttt{flex-direction: column} ao \texttt{ul}.
\end{frame}

\begin{frame}[fragile]{A grade de produtos}
	\begin{minted}{css}
.grade {
  display: grid;
  grid-template-columns: repeat(auto-fit, minmax(240px, 1fr));
  gap: 20px;
}

img { max-width: 100%; height: auto; }
  \end{minted}

	\vspace{0.2cm}
	Ligue o desenho da grade no inspetor e conte as colunas. Estreite a janela até
	a grade passar para duas colunas, anote a largura e confira a conta.

	\vspace{0.2cm}
	Depois troque \texttt{240px} por \texttt{320px} e observe onde ela muda.
\end{frame}

\begin{frame}[fragile]{O cartão por dentro}
	Descrições de tamanhos diferentes deixam o preço em altura diferente em cada
	cartão.

	\begin{minted}{css}
.produto {
  display: flex;
  flex-direction: column;
  height: 100%;
}

.produto .preco {
  margin-top: auto;  /* absorve a sobra vertical */
}
  \end{minted}

	\vspace{0.2cm}
	\texttt{margin: auto} dentro de um contêiner Flex absorve o espaço livre
	daquele lado. É a forma mais curta de empurrar um item para a ponta do eixo.
\end{frame}

\begin{frame}[fragile]{Media queries e o teste final}
	\begin{minted}{css}
header { flex-direction: column; }

@media (min-width: 768px) {
  header { flex-direction: row; }
}
  \end{minted}

	\vspace{0.2cm}
	Em cada uma das três páginas:

	\begin{enumerate}
		\item \texttt{Ctrl+Shift+M} para ligar o modo responsivo;
		\item largura 320, e rolar a página de cima a baixo;
		\item achar no inspetor o elemento desenhado além da borda direita;
		\item repetir em 1200px e confirmar que o layout largo continua correto.
	\end{enumerate}
\end{frame}

\begin{frame}{Exercícios}
	\begin{block}{Entrega}
		Valem nota, no item \emph{Exercícios em sala}. O \textbf{prazo} está na
		atividade correspondente na \textbf{UFPR Virtual}, e a entrega é por
		\texttt{push} no GitLab.
	\end{block}

	\vspace{0.4cm}
	\begin{itemize}
		\item enunciado no \texttt{README.md} do repositório-modelo, com o endereço
		      na UFPR Virtual;
		\item individual ou em dupla, com \emph{fork} dentro do grupo;
		\item atividade avaliativa: sem IA generativa para produzir o código;
		\item o que sair daqui é a \textbf{Entrega 1} do trabalho prático.
	\end{itemize}

	\vspace{0.3cm}
	Para treinar fora da entrega, em português: Flexbox Froggy e Grid Garden.
\end{frame}

\begin{frame}{Resumo: Flexbox e Grid}
	\begin{itemize}
		\item os dois modelos entram quando o \textbf{pai} declara o
		      \texttt{display}, e valem para os filhos diretos;
		\item propriedade de contêiner escrita no item não faz nada e não gera erro;
		\item \texttt{justify-content} no eixo principal, \texttt{align-items} no
		      transversal, e \texttt{flex-direction} troca o sentido dos dois;
		\item \texttt{flex-wrap: nowrap} é o padrão: sem \texttt{wrap} os itens
		      encolhem em vez de quebrar a linha;
		\item \texttt{repeat(auto-fit, minmax(240px, 1fr))} muda o número de colunas
		      sozinho, sem \emph{media query};
		\item Grid para o esqueleto e para grades, Flexbox para uma linha ou uma
		      coluna de itens.
	\end{itemize}
\end{frame}

\begin{frame}{Resumo: responsividade}
	\begin{itemize}
		\item sem a meta viewport, o celular desenha numa janela fingida de 980px e
		      nenhuma \emph{media query} estreita entra em ação;
		\item em \emph{mobile-first}, as \emph{media queries} usam
		      \texttt{min-width} e vão ao final da folha, porque não alteram
		      especificidade;
		\item o ponto de quebra sai de onde o layout fica ruim, e não de uma lista
		      de aparelhos;
		\item a 320px de largura não pode haver rolagem horizontal (WCAG 1.4.10);
		\item Bootstrap: 12 colunas e utilitárias, ao custo de 227\,KB e de HTML
		      cheio de classes de apresentação. Aqui, conteúdo de leitura.
	\end{itemize}
\end{frame}

\end{document}
